% Generated draft: requires mathematical and source review.
\documentclass[11pt]{article}
\usepackage[margin=25mm]{geometry}
\usepackage{fontspec}
\setmainfont{Latin Modern Roman}
\usepackage{amsmath,amssymb,mathtools,unicode-math}
\setmathfont{Latin Modern Math}
\usepackage{enumitem,needspace,xurl,hyperref}
\setlength{\parindent}{0pt}
\setlength{\parskip}{5pt}
\setlength{\emergencystretch}{3em}
\newcommand{\N}{\mathbb{N}}
\newcommand{\Z}{\mathbb{Z}}
\newcommand{\Q}{\mathbb{Q}}
\newcommand{\R}{\mathbb{R}}
\newcommand{\C}{\mathbb{C}}
\newcommand{\F}{\mathbb{F}}
\newcommand{\A}{\mathbb{A}}
\newcommand{\PP}{\mathbb{P}}
\newcommand{\E}{\mathbb{E}}
\newcommand{\eps}{\varepsilon}
\newcommand{\dd}{\,\mathrm{d}}
\begin{document}
\section*{UnsolvedMath ID: 20002098}
\textbf{Source problem code:} catalogue entry
\subsection{Definitions and mathematical statement}
The complete numbered Question 1.40 says only that Mitsumatsu discussed linking numbers of vector-field orbits in relation to Reeb fields and foliations. It contains no determinate conjecture, hypotheses, or requested conclusion. The aid explicitly does not reconstruct a lost question. The following is its contextual three-dimensional model, not an asserted recovery of the original intended statement.

Let \(M\) be a closed oriented three-manifold. For exact two-forms define the helicity pairing
\[
 \mathcal B(\dd\alpha,\dd\beta)=\int_M\alpha\wedge\dd\beta.
\]
The integral is independent of the primitive choices. If \(\alpha\wedge\dd\alpha\) is a positive contact volume, the restriction to the real line generated by \(\dd\alpha\) is positive definite. For a foliation-defining one-form \(\theta\) with \(\theta\wedge\dd\theta=0\), the line generated by \(\dd\theta\) is null for this quadratic pairing (or is the zero space if \(\dd\theta=0\)). No positive-definiteness claim is made for the whole space of exact two-forms.

The Reeb field \(R_\alpha\), defined by \(\alpha(R_\alpha)=1\) and \(\iota_{R_\alpha}\dd\alpha=0\), has exact flux with respect to \(\mu=\alpha\wedge\dd\alpha\):
\[
 \iota_{R_\alpha}\mu=\dd\alpha,\qquad
 \mathcal H(R_\alpha;\mu)=\int_M\alpha\wedge\dd\alpha.
\]
The volume convention is not normalized to a probability measure. If an additional adapted open book is given, positive Reeb crossing of pages gives positive oriented intersection with a fixed page. Calling this a knot linking number requires that the page and orbit satisfy the requisite homological conditions; for example on an oriented integral homology three-sphere a page bounding the binding computes its linking with a disjoint orbit. No such open book or homological setting is part of the original one-sentence record.

For its explicit family, fix \(a,b>0\) and use the normalization
\[
 E(a,b)=\{(z_1,z_2)\in\C^2:\pi|z_1|^2/a+\pi|z_2|^2/b=1\},
 \qquad \alpha=\tfrac12\sum_{j=1}^2(x_j\dd y_j-y_j\dd x_j)|_{E(a,b)}.
\]
Let \(B_1=\{z_2=0\}\) and \(B_2=\{z_1=0\}\), oriented by the Reeb flow. In this convention their minimal periods are \(T(B_1)=a\), \(T(B_2)=b\), and contact helicity is \(\mathcal H=ab\). The aid reports for every other periodic Reeb orbit \(\delta\), disjoint from the chosen axis,
\[
 \mathcal H\,\frac{\operatorname{lk}(B_i,\delta)}{T(\delta)}=T(B_i).
\]
For a multiply covered orbit both the period and linking are counted with its covering multiplicity; for a knot use its minimal period.

The pages around \(B_1\) are given by fixed \(\arg z_2\), and around \(B_2\) by fixed \(\arg z_1\). For any invariant probability measure supported away from the corresponding axis, define \(\rho_i\) as the mean positive page rate in turns per unit Reeb time. The aid's measure version is
\[
 \mathcal H\rho_i=T(B_i),\qquad \rho_1=1/b,\quad\rho_2=1/a.
\]
Thus the identities compare analytic helicity, periodic linking, and mean page crossing for this family. They neither supply a general linking theorem for every Reeb flow nor answer an unspecified foliation question.
\subsection{Short English statement}
The source contains a discussion note rather than a mathematical question. The aid develops a helicity model and an exact ellipsoid identity linking axis periods to orbit linking and page-crossing rates. Open-book and homological assumptions are kept explicit rather than inferred from the brief source sentence.
\subsection{Sources}
\begin{itemize}
\item Local UnsolvedMath registry, permanent ID 20002098, exact source code catalogue entry; canonical \texttt{aim-geometry-notes.json}, \texttt{notes[435]}, tag \texttt{question}; AIM Workshop Problem Lists, Geometry, \emph{Holomorphic curves in contact geometry}, Question 1.40. Created and updated 14 August 2026.
\item Source URL \url{https://aimath.org/WWN/contactgeom2/contactgeom2.pdf}. The entire supplied source attributes a discussion to Mitsumatsu and gives no determinate conjecture. The PDF was not accessed.
\item No complete authoritative reference is supplied for the helicity pairing, adapted-open-book argument, or standard-ellipsoid identity in the aid.
\item Attempt 1: \texttt{context\_only}, problem status at run \texttt{unknown}, rigor label \texttt{proved}, novelty confidence \texttt{low}; classification \texttt{OPEN-TRIAGE}, suggested difficulty L2, registry status \texttt{open}. Candidate: explicit helicity--linking--period and invariant-measure identities. These are contextual source claims, not a solved reconstructed question.
\item Ulam AI / UnsolvedMath Contributors; CC BY 4.0. Machine-generated research aid, not peer reviewed; independently verify before citation.
\end{itemize}
\end{document}
